\magicSubsection{Swizzle Period}{sec:gSwizzlePeriod}

A \lighttt{sm\_90} swizzle function (all we support for now) works by XOR-ing certain bits in the pointer with other bits in the pointer.
The common 128-byte swizzle function works like this (using Verilog syntax)
\begin{verbatim}
        address[6:4] ^= address[9:7]
\end{verbatim}

In general, once a pointer is swizzled, it cannot be adjusted with pointer arithmetic anymore.
The exception: there exists a minimal swizzle period $P$ such that adjustments that are multiples of $P$ bytes can be done after the swizzling.
i.e., the swizzle pattern of the swizzle function $S$ repeats every $P$ bytes:
\begin{equation*}
    \forall x, y \in \mathbb{Z}. S(Px + y) = Px + S(y)
\end{equation*}
The 128-byte swizzle function, for example, has a swizzle period of 1024 bytes.

Non-swizzled memory may be viewed as having an identity swizzle function with period $P = 1$.
